% ============================================================================
% APPENDIX MM: LIT-MARECHAL - Honesty & Corporate Culture
% ============================================================================
% Category: Literature Integration (LIT-)
% Language: English
% EBF Integration: Corporate culture, identity priming, and honesty research
% ============================================================================

\chapter*{Appendix UNMAPPED_MM: Michel Maréchal Research Integration}
\addcontentsline{toc}{chapter}{Appendix UNMAPPED_MM: LIT-MARECHAL - Honesty \& Corporate Culture}
\label{app:marechal}

% ============================================================================
% HEADER BLOCK
% ============================================================================
\begin{tcolorbox}[colback=white,colframe=black!75!white,title=Appendix Metadata]
\textbf{Appendix:} MM -- LIT-MARECHAL \\
\textbf{Category:} Literature Integration (LIT-) \\
\textbf{Version:} 1.0 (January 2026) \\
\textbf{Dependencies:} CORE-WHEN, CONTEXT-INSTITUTION, LIT-FEHR, LIT-FISCHBACHER \\
\textbf{Status:} Complete
\end{tcolorbox}

% ============================================================================
% CROSS-REFERENCE MAP
% ============================================================================
\begin{tcolorbox}[colback=blue!5!white,colframe=blue!75!black,title=Cross-Reference Map]
\textbf{Builds On:}
\begin{itemize}
    \item Appendix K (LIT-FEHR): Social preferences foundations
    \item Appendix UNMAPPED_FI (LIT-FISCHBACHER): Experimental methodology
    \item Appendix UNMAPPED_V (CORE-WHEN): Context dependence $\Psi$
    \item Appendix AH (CONTEXT-INSTITUTION): Institutional effects
\end{itemize}

\textbf{Extends To:}
\begin{itemize}
    \item Appendix AD (DOMAIN-FINANCE): Financial professional behavior
    \item Appendix AG (DOMAIN-HR): Workplace dynamics
    \item Appendix AO (PREDICT-BEHAVIOR): Behavioral prediction
\end{itemize}

\textbf{Methods Connection:}
\begin{itemize}
    \item Appendix E (METHOD-E): Field experimental methodology
    \item Appendix UNMAPPED_LLM (METHOD-LLMMC): Cross-cultural estimation
\end{itemize}
\end{tcolorbox}

% ============================================================================
% CHAPTER LINKAGE
% ============================================================================
\begin{tcolorbox}[colback=green!5!white,colframe=green!75!black,title=Chapter Linkage]
\textbf{Primary:} Chapter 10 (Context Effects) \\
\textbf{Secondary:} Chapter 5 (Social Preferences), Chapter 14 (Organizational Applications)
\end{tcolorbox}

% ============================================================================
% ABSTRACT
% ============================================================================
\begin{abstract}
This appendix integrates Michel Maréchal's research on honesty, corporate culture, and cross-cultural behavior into the Evidence-Based Framework. Maréchal's work demonstrates how institutional context---particularly corporate culture and professional identity---shapes ethical behavior, with profound implications for organizational design and policy. His global studies reveal both universal patterns and cultural variation in civic honesty. The research provides critical parameters for the context function $\Psi$ and establishes predictive links between laboratory behavior and real-world outcomes.
\end{abstract}

% ============================================================================
% QUICK REFERENCE
% ============================================================================
\begin{tcolorbox}[colback=gray!10!white,colframe=gray!50!black,title=Quick Reference]
Key terms defined in this appendix (see also Appendix G for complete Glossary):
\begin{description}
    \item[Identity Priming] Activating a specific social identity to study its effects on behavior
    \item[Corporate Culture] Shared norms, values, and practices within an organization
    \item[Civic Honesty] Propensity to return lost property to its rightful owner
    \item[Moral Universalism] Treating in-group and out-group members equally in moral decisions
\end{description}
\end{tcolorbox}

% ============================================================================
% FUNDAMENTAL QUESTION
% ============================================================================
% -----------------------------------------------------------------------------
% APPENDIX SCOPE BOX
% -----------------------------------------------------------------------------

\begin{tcolorbox}[
    colback=orange!5!white,
    colframe=orange!75!black,
    title={\textbf{Appendix Scope: Ziel / In-Scope / Out-of-Scope / Constraints}},
    fonttitle=\bfseries
]

\textbf{Ziel (Objective):}
\begin{itemize}[nosep]
    \item How does institutional context---particularly corporate culture and professional identity---shape ethical behavior, and can these effects be measured, predicted, and modified?
\end{itemize}

\textbf{In-Scope:}
\begin{itemize}[nosep]
    \item Fundamental Question
    \item Theoretical Framework
    \item Core Axioms
    \item Key Empirical Results
    \item EBF Integration
\end{itemize}

\textbf{Out-of-Scope (delegiert an andere Appendices):}
\begin{itemize}[nosep]
    \item LIT-FEHR $\rightarrow$ Appendix K
    \item LIT-FISCHBACHER $\rightarrow$ Appendix UNMAPPED_FI
    \item CORE-WHEN $\rightarrow$ Appendix UNMAPPED_V
    \item CONTEXT-INSTITUTION $\rightarrow$ Appendix AH
    \item DOMAIN-FINANCE $\rightarrow$ Appendix AD
\end{itemize}

\textbf{Constraints (Anwendungsgrenzen):}
\begin{itemize}[nosep]
    \item [TODO: Define when this appendix does NOT apply]
    \item [TODO: Define required prerequisites]
    \item [TODO: Define known limitations]
\end{itemize}

\textbf{Lieferobjekte (Deliverables):}
\begin{itemize}[nosep]
    \item 4 Table(s)
    \item 10 Equation(s)
    \item 6 Axiom(s)
    \item Worked Example(s)
\end{itemize}

\end{tcolorbox}


\section{Fundamental Question}

\begin{tcolorbox}[colback=yellow!10!white,colframe=orange!75!black,title=Central Research Question]
\textbf{How does institutional context---particularly corporate culture and professional identity---shape ethical behavior, and can these effects be measured, predicted, and modified?}
\end{tcolorbox}

This question addresses CORE-WHEN ($\Psi$): When and how does context systematically alter behavior?

% ============================================================================
% THEORETICAL FRAMEWORK
% ============================================================================
\section{Theoretical Framework}

\subsection{The Context-Honesty Nexus}

Maréchal's research establishes that honesty is not a fixed trait but emerges from the interaction of individual dispositions and contextual factors:

\begin{equation}
H_i(t) = \bar{H}_i + \psi_{culture} \cdot C_{org} + \psi_{identity} \cdot I_{active}(t) + \epsilon_i(t)
\end{equation}

Where:
\begin{itemize}
    \item $H_i(t)$ = Individual $i$'s honesty at time $t$
    \item $\bar{H}_i$ = Baseline honesty disposition
    \item $C_{org}$ = Organizational culture parameter
    \item $I_{active}(t)$ = Currently activated identity
    \item $\psi_{culture}, \psi_{identity}$ = Context sensitivity parameters
\end{itemize}

\subsection{The Corporate Culture Channel}

The seminal Cohn, Fehr \& Maréchal (2014) Nature paper demonstrates that professional identity in banking increases dishonest behavior:

\begin{equation}
\Delta H_{banker} = H_{control} - H_{primed} \approx 16\% \text{ (increased cheating)}
\end{equation}

This effect operates through identity salience:
\begin{equation}
I_{banker}^{active} \xrightarrow{\text{priming}} \text{Competitive norms} \xrightarrow{} \downarrow H
\end{equation}

% ============================================================================
% AXIOMS
% ============================================================================
\section{Core Axioms}

\begin{tcolorbox}[colback=red!5!white,colframe=red!75!black,title=Axiom UNMAPPED_MM-1: Corporate Culture Effect]
\textbf{Corporate culture systematically shapes individual ethical behavior.}

Formally:
\begin{equation}
\frac{\partial H_i}{\partial C_{org}} = \psi_{culture} > 0
\end{equation}

Employees in ethical cultures behave more honestly; those in competitive/materialistic cultures behave less honestly.

\textit{Evidence:} Cohn, Fehr \& Maréchal (2014), $N=208$ bankers, Nature.
\end{tcolorbox}

\begin{tcolorbox}[colback=red!5!white,colframe=red!75!black,title=Axiom UNMAPPED_MM-2: Identity Priming Effect]
\textbf{Activating professional identity alters behavior consistent with that identity's norms.}

\begin{equation}
H_i | I_{banker}^{active} < H_i | I_{private}^{active}
\end{equation}

Banking identity priming increases dishonesty by 16 percentage points relative to control.

\textit{Evidence:} Cohn, Fehr \& Maréchal (2014), banker vs. non-finance professional comparison.
\end{tcolorbox}

\begin{tcolorbox}[colback=red!5!white,colframe=red!75!black,title=Axiom UNMAPPED_MM-3: Civic Honesty Universality with Variation]
\textbf{Civic honesty is universal but varies systematically with institutional quality and cultural values.}

\begin{equation}
P(\text{wallet returned}) = f(\text{money in wallet}, \text{country}, \text{institutional quality})
\end{equation}

Key finding: More money $\rightarrow$ higher return rates (opposite of pure self-interest prediction).

\textit{Evidence:} Cohn, Maréchal et al. (2019), 40 countries, 17,000+ wallets, Science.
\end{tcolorbox}

\begin{tcolorbox}[colback=red!5!white,colframe=red!75!black,title=Axiom UNMAPPED_MM-4: Neural Basis of Honesty]
\textbf{Honesty can be causally increased through right prefrontal cortex stimulation.}

\begin{equation}
H | TMS_{rPFC} > H | Sham \quad (\Delta \approx 20\%)
\end{equation}

Transcranial magnetic stimulation increases honest behavior by enhancing cognitive control.

\textit{Evidence:} Maréchal, Cohn \& Ugazio (2017), PNAS.
\end{tcolorbox}

\begin{tcolorbox}[colback=red!5!white,colframe=red!75!black,title=Axiom UNMAPPED_MM-5: Lab-Field Predictive Validity]
\textbf{Laboratory measures of cheating predict real-world misconduct.}

\begin{equation}
P(\text{school misconduct}) = \alpha + \beta \cdot \text{Lab Cheating} + \epsilon, \quad \beta > 0
\end{equation}

Lab cheating tasks have predictive validity for actual behavioral outcomes.

\textit{Evidence:} Cohn \& Maréchal (2018), Economic Journal.
\end{tcolorbox}

\begin{tcolorbox}[colback=red!5!white,colframe=red!75!black,title=Axiom UNMAPPED_MM-6: Honesty-Cognition Link]
\textbf{Honest individuals tend to use less strategic reasoning.}

\begin{equation}
\text{Theory of Mind Depth} | H_{high} < \text{ToM Depth} | H_{low}
\end{equation}

Honest people are less likely to engage in sophisticated strategic thinking about others' beliefs.

\textit{Evidence:} Maréchal, Cohn, Ugazio \& Ruff (2019), PNAS.
\end{tcolorbox}

% ============================================================================
% KEY EMPIRICAL RESULTS
% ============================================================================
\section{Key Empirical Results}

\subsection{The Banking Culture Study (Nature 2014)}

\begin{table}[h]
\centering
\begin{tabular}{lcc}
\hline
\textbf{Condition} & \textbf{Reported Wins} & \textbf{Excess over Fair} \\
\hline
Bankers - Control & 51.6\% & 1.6\% \\
Bankers - Identity Prime & 58.2\% & 8.2\% \\
Non-finance - Identity Prime & 51.8\% & 1.8\% \\
\hline
\end{tabular}
\caption{Coin flip reporting in banking identity study}
\end{table}

\textbf{Key insight:} Banking identity---not financial incentives per se---triggers the effect.

\subsection{Global Civic Honesty (Science 2019)}

\begin{table}[h]
\centering
\begin{tabular}{lccc}
\hline
\textbf{Country Type} & \textbf{No Money} & \textbf{With Money} & \textbf{Difference} \\
\hline
High GDP/Institutions & 46\% & 61\% & +15pp \\
Low GDP/Institutions & 32\% & 45\% & +13pp \\
\hline
\end{tabular}
\caption{Wallet return rates by country type and money condition}
\end{table}

\textbf{Key insight:} Against economic intuition, wallets with more money are returned more often.

\subsection{Cross-Cultural Preference Heterogeneity}

From Falk et al. (2018, QJE), with Maréchal as co-author:

\begin{itemize}
    \item 80,000 respondents across 76 countries
    \item 6 preference dimensions: risk, time, altruism, reciprocity, trust, negative reciprocity
    \item Cultural clusters explain significant variance
    \item Within-country variation exceeds between-country variation
\end{itemize}

% ============================================================================
% EBF INTEGRATION
% ============================================================================
\section{EBF Integration}

\subsection{Context Function Specification}

Maréchal's research provides critical parameters for the context function $\Psi$:

\begin{equation}
\Psi_{honesty} = \begin{pmatrix}
\psi_{corporate} \\
\psi_{identity} \\
\psi_{culture} \\
\psi_{institution}
\end{pmatrix}
\end{equation}

\begin{table}[h]
\centering
\begin{tabular}{lccl}
\hline
\textbf{Parameter} & \textbf{Estimate} & \textbf{SE} & \textbf{Source} \\
\hline
$\psi_{corporate}$ & 0.16 & 0.04 & Nature 2014 \\
$\psi_{identity}$ & 0.12 & 0.03 & Multiple \\
$\psi_{culture}$ & 0.08 & 0.02 & Science 2019 \\
$\psi_{institution}$ & 0.15 & 0.04 & Science 2019 \\
\hline
\end{tabular}
\caption{Context sensitivity parameters from Maréchal research}
\end{table}

\subsection{10C Framework Mapping}

\begin{table}[h]
\centering
\begin{tabular}{lll}
\hline
\textbf{CORE} & \textbf{Maréchal Contribution} & \textbf{Key Paper} \\
\hline
WHO (AAA) & Professional identity effects & Nature 2014 \\
WHAT (C) & Moral preferences as utility component & Ann Rev Econ 2020 \\
HOW (B) & Norms as coordination device & Nat Hum Behav 2022 \\
WHEN (V) & Context-dependent honesty & All work \\
WHERE (BBB) & Global parameters & QJE 2018, Science 2019 \\
AWARE (AU) & Strategic reasoning \& honesty & PNAS 2019 \\
\hline
\end{tabular}
\caption{10C Framework mapping for Maréchal research}
\end{table}

% ============================================================================
% WORKED EXAMPLE
% ============================================================================
\section{Worked Example: Designing Ethical Banking Culture}

\begin{tcolorbox}[colback=blue!5!white,colframe=blue!75!black,title=Application: Post-Crisis Bank Culture Reform]
\textbf{Context:} A major bank seeks to improve ethical culture after misconduct scandals.

\textbf{Step 1: Diagnose Current State}
Using Maréchal parameters:
\begin{itemize}
    \item Current identity priming: Banking identity strongly associated with competition
    \item $\psi_{corporate} \approx -0.20$ (negative effect on honesty)
\end{itemize}

\textbf{Step 2: Design Intervention}
\begin{enumerate}
    \item \textbf{Re-frame Professional Identity:} Emphasize stewardship, fiduciary duty
    \item \textbf{Modify Cultural Symbols:} Replace competition-focused messaging
    \item \textbf{Structural Changes:} Align incentives with long-term client outcomes
\end{enumerate}

\textbf{Step 3: Predict Effect}
Target: Shift $\psi_{corporate}$ from $-0.20$ to $+0.05$

Expected change in honest behavior: $\Delta H \approx 0.25 \times 0.16 = 4\%$ reduction in misconduct

\textbf{Step 4: Measurement}
\begin{itemize}
    \item Pre/post identity priming experiments (internal)
    \item Track misconduct metrics
    \item Employee surveys on ethical climate
\end{itemize}
\end{tcolorbox}

% ============================================================================
% IMPLICATIONS
% ============================================================================
\section{Implications for Practice}

\subsection{For Organizations}

\begin{enumerate}
    \item \textbf{Culture Audits:} Measure implicit norms through priming studies
    \item \textbf{Identity Management:} Carefully curate professional identity messaging
    \item \textbf{Hiring:} Lab-based honesty measures have predictive validity
    \item \textbf{Training:} Focus on ethical identity rather than rule compliance
\end{enumerate}

\subsection{For Policy}

\begin{enumerate}
    \item \textbf{Regulatory Design:} Account for culture, not just incentives
    \item \textbf{Institutional Quality:} Improves civic honesty baseline
    \item \textbf{Cross-Cultural Programs:} Universal mechanisms exist but calibration needed
\end{enumerate}

% ============================================================================
% SUMMARY
% ============================================================================
\section{Summary}

\begin{tcolorbox}[colback=gray!10!white,colframe=gray!50!black,title=Key Takeaways]
\begin{enumerate}
    \item \textbf{Context Matters:} Corporate culture and identity priming causally affect honesty
    \item \textbf{Universal Mechanisms:} Civic honesty follows predictable patterns globally
    \item \textbf{Predictive Validity:} Lab measures predict real-world behavior
    \item \textbf{Neural Basis:} Honesty can be causally manipulated via brain stimulation
    \item \textbf{Design Implications:} Organizational design must account for cultural effects
\end{enumerate}

\textbf{Key Insight}: Ethical behavior emerges from context $\times$ disposition interaction. Design the context.
\end{tcolorbox}

% ============================================================================
% GLOSSARY
% ============================================================================
\section{Glossary}

Key terms defined in this appendix (see also Appendix G for complete Glossary):

\begin{description}
    \item[Identity Priming] Experimental technique activating specific social identities
    \item[Corporate Culture] Shared organizational norms affecting individual behavior
    \item[Civic Honesty] Tendency to return lost property (wallet paradigm)
    \item[Moral Universalism] Applying same moral standards to in-group and out-group
    \item[Strategic Reasoning] Thinking about others' beliefs and intentions
    \item[rPFC Stimulation] Non-invasive brain stimulation of right prefrontal cortex
\end{description}

% ============================================================================
% CRITICAL FOUNDATIONS
% ============================================================================
\section{Critical Foundations}

\begin{enumerate}
    \item \textbf{Akerlof \& Kranton (2000):} Identity economics framework
    \item \textbf{Ariely (2012):} Dishonesty as pervasive but malleable
    \item \textbf{Gneezy et al. (2005):} Deception games methodology
    \item \textbf{Fehr \& Fischbacher (2003):} Social preferences foundations
    \item \textbf{Henrich et al. (2010):} Cross-cultural behavioral variation
\end{enumerate}

% ============================================================================
% OPEN ISSUES
% ============================================================================
\section{Open Issues}

\begin{enumerate}
    \item \textbf{Long-term Effects:} Do culture interventions persist?
    \item \textbf{Selection vs. Socialization:} What proportion of culture effect is hiring vs. adaptation?
    \item \textbf{Optimal Regulation:} How should regulators target culture vs. incentives?
    \item \textbf{Measurement:} Can we develop real-time culture diagnostics?
    \item \textbf{Generalization:} Do banking findings extend to other professions?
\end{enumerate}

% ============================================================================
% REFERENCES
% ============================================================================
\section{References}

See Master Bibliography (bcm2\_master\_references.bib) for complete citations.

\nocite{bcm_master}

\textbf{Primary Maréchal Sources} (16 papers integrated):

\begin{itemize}
    \item Cohn, A., Fehr, E., \& Maréchal, M.A. (2014). Business Culture and Dishonesty in the Banking Industry. \textit{Nature}, 516, 86--89.
    \item Cohn, A., Maréchal, M.A., Tannenbaum, D., \& Zünd, C.L. (2019). Civic Honesty Around the Globe. \textit{Science}, 365, 70--73.
    \item Falk, A., et al. (2018). Global Evidence on Economic Preferences. \textit{QJE}, 133(4), 1645--1692.
    \item Maréchal, M.A., Cohn, A., \& Ugazio, G. (2017). Increasing Honesty with Noninvasive Brain Stimulation. \textit{PNAS}, 114(17).
    \item Kube, S., Maréchal, M.A., \& Puppe, C. (2012). The Currency of Reciprocity. \textit{AER}, 102(4).
    \item Kube, S., Maréchal, M.A., \& Puppe, C. (2013). Do Wage Cuts Damage Work Morale? \textit{JEEA}, 11(4).
    \item Cohn, A., \& Maréchal, M.A. (2018). Lab Cheating Predicts School Misconduct. \textit{EJ}, 128.
    \item Cohn, A., Fehr, E., \& Maréchal, M.A. (2015). Countercyclical Risk Aversion. \textit{AER}, 105(2).
    \item Fehr, E., Schurtenberger, I., \& Maréchal, M.A. (2022). Normative Foundations of Cooperation. \textit{Nat Hum Behav}.
    \item Cohn, A., Maréchal, M.A., Schneider, F., \& Weber, R.A. (2022). Moral Universalism. \textit{JPE}, 130(7).
\end{itemize}

\end{document}
